\bookchapter{제38장: AI 프로그래밍 5가지 사고 원칙}

\chapterintro{이 책에서 단 한 장만 기억할 수 있다면 그건 반드시 이 장이어야 해요. 5가지 사고 원칙은 책 전체의 심장이고, 뒤의 모든 장——컨텍스트, 요구사항, 대화, 품질——은 이 다섯 가지 원칙을 각기 다른 상황에서 펼쳐 놓은 것이니까요. 걱정 마세요. 하나하나 어렵지 않지만, 하나하나 모두 몇 개의 프로젝트를 통해 '몸에 새기는' 가치가 있어요.}

이야기의 출발점이 재미있어요. Andrej Karpathy가 자기 프로젝트를 오픈소스로 공개할 때 단 70줄짜리 CLAUDE.md 파일을 곁들였는데, 거기에는 AI와 협업하는 네 가지 핵심 원칙이 정리되어 있었어요. 이 파일은 GitHub에서 70,000개가 넘는 Star를 받으며 전 세계 AI 프로그래밍 실천가들이 가장 널리 인용하는 기준이 되었죠. 이 네 가지 원칙은 텅 빈 구호가 아니라 수많은 실전에서 추려낸, 실행 가능한 사고 습관이에요. 이 책은 여기에 다섯 번째 원칙——반복 검증——을 더해 완전한 5대 원칙 체계를 구성해요.

\begin{tcbinfo}{70줄짜리 파일은 왜 이렇게 가치가 높을까요?}
Star 70,000개면 어떤 개념일까요? GitHub의 프로젝트 대부분은 평생 이 숫자에 도달하지도 못해요. 70줄짜리 텍스트 파일이 그걸 받았다는 건 딱 한 가지를 증명해요. AI 시대에 가장 희소한 것은 코드가 아니라 '명확하게 생각하는' 능력이라는 거죠. 이 70줄은 바로 '명확한 사고'의 농축판입니다.

\end{tcbinfo}

\booksection{원칙 1: Think Before Coding(먼저 생각하고, 그다음 코딩하기)}

다섯 가지 원칙 중 가장 직관에 어긋나는 원칙이에요. AI가 코드를 짜 줄 수 있는데 왜 '먼저 생각'해야 할까요? 답은 간단해요. AI의 결과물 품질은 여러분의 입력 품질에 비례해요. 자기도 뭘 만들지 명확히 생각하지 못한 상태라면, AI가 내놓는 답은 끝없는 '고치고 또 고치기' 순환에 빠지게 만들 뿐이죠. 흔히 말하는 'AI가 같이 길을 잃어 주는' 상황이에요.

'먼저 생각하기'란 AI에게 지시를 내리기 전에 세 가지 질문에 답해야 한다는 뜻이에요. 이 세 가지는 하나 빠지면 안 되는 '생각의 삼각형'을 이뤄요.

\begin{enumerate}
  \item 나는 무슨 문제를 해결하려는 걸까?(목표)
  \item 내가 기대하는 결과는 어떤 모습일까?(검수 기준)
  \item 지켜야 할 제약 조건은 무엇일까?(경계)
\end{enumerate}

\begin{center}
\begin{tabularx}{\linewidth}{|X|X|X|}
\hline
\textbf{생각 항목} & \textbf{모호한 입력(나쁨)} & \textbf{명확한 입력(좋음)} \tabularnewline
\hline
문제 정의 & 웹사이트 하나 만들어 줘 & AI 소도구 3개의 입구 카드를 보여 주는 개인 도구 사이트를 만들어 줘 \tabularnewline
기대 결과 & 좀 예쁘게 & 반응형 레이아웃, 다크 테마, 카드에 호버 애니메이션, 모바일에서도 사용 가능 \tabularnewline
제약 조건 & (미설정) & 프레임워크 없이 순수 HTML+CSS+JS, 단일 파일 배포, 백엔드 불필요 \tabularnewline
\hline
\end{tabularx}
\end{center}

\begin{center}\small\color{gray}같은 요구에 대한 모호한 입력과 명확한 입력 비교\end{center}

\begin{tcbwarn}{주의: 먼저 생각하기 $\neq$ 완벽한 문서부터 쓰기}
'먼저 생각하기'는 '완벽한 PRD를 쓰고 시작하기'와 같지 않아요. 빠르지만 질서 있는 사고 과정이에요. 30초만 들여 위의 세 가지 질문을 명확히 하는 편이, 대화 속에서 30분을 고치고 또 고치는 것보다 훨씬 효율적이에요. 작은 과제라면 세 질문은 아주 짧아도 되고, 큰 프로젝트라면 이를 완전한 요구사항 문서로 확장해도 돼요(자세한 내용은 제40장).

\end{tcbwarn}

\booksection{원칙 2: Simplicity First(단순함이 최우선)}

AI 프로그래밍의 가장 큰 함정은 '못 만드는 것'이 아니라 '너무 많이 만드는 것'이에요. AI에게 '사용자 관리 시스템을 만들어 줘'라고 하면 인증, 권한, 역할, 감사 로그, 데이터 내보내기 등 온갖 기능을 한꺼번에 만들어 낼 수 있어요. 하지만 지금 여러분에게 필요한 건 간단한 로그인 페이지 하나일 수도 있죠. 과도한 생성이야말로 AI 코드가 부풀어 오르는 주된 원인이에요.

'단순함이 최우선'은 지시를 내릴 때마다 '최소한으로 동작하는' 원칙을 지키라고 요구해요. 가장 적은 코드로 핵심 기능을 구현하고, 가장 단순한 방안으로 문제를 해결하고, 가장 적은 기능으로 아이디어를 검증하는 거죠. Karpathy는 이에 대한 통찰력 있는 비유를 남겼어요. '할인 가격 하나 계산하려고 전자상거래 엔진이 필요하지는 않다.'

실전에서 '단순함이 최우선'은 몇 가지 구체적인 습관을 뜻해요.

\begin{itemize}
  \item 한 번에 AI에게 기능 하나만 구현하게 하고, 시스템 전체를 한꺼번에 설명하지 않아요.
  \item 복잡한 의존성을 도입하기보다 표준 라이브러리와 단순한 방안을 먼저 선택해요.
  \item 50줄로 해결할 수 있는 문제라면 AI에게 500줄을 쓰게 하지 않아요.
  \item AI의 결과물이 기대 범위를 벗어나면 용기 내어 '필요 없는 부분은 잘라 내라고' 말하세요.
\end{itemize}

\begin{center}
\begin{tabularx}{\linewidth}{|X|X|X|}
\hline
\textbf{상황} & \textbf{과도한 방안(피할 것)} & \textbf{단순한 방안(추천)} \tabularnewline
\hline
데이터를 저장해야 할 때 & MySQL 데이터베이스 + ORM 구축 & 일단 JSON 파일이나 SQLite로 시작 \tabularnewline
사용자 로그인이 필요할 때 & OAuth2.0 + JWT + 역할 기반 권한 & 일단 간단한 비밀번호 검증부터 \tabularnewline
프론트엔드 프레임워크가 필요할 때 & React + Redux + Router & 일단 순수 HTML + 네이티브 JS로 시도 \tabularnewline
배포해야 할 때 & Docker + CI/CD + K8s & 일단 로컬에서 돌아가는지부터 확인 \tabularnewline
\hline
\end{tabularx}
\end{center}

\begin{center}\small\color{gray}자주 있는 상황에서의 과도한 방안과 단순한 방안 비교\end{center}

\booksection{원칙 3: Surgical Changes(정밀 수정)}

이미 있는 프로젝트를 반복 개선할 때는 매번 수정을 외과 수술처럼 정밀하게 해야 해요. 움직여야 할 곳만 움직이고, 건드리지 말아야 할 곳은 건드리지 않는 거죠. 이건 AI 프로그래밍에서 특히 중요해요. AI에는 '지나가다 한번 손봐 주는' 성향이 있거든요. 버그 하나 고치라고 했는데 파일 절반을 리팩터링해 버릴 수도 있어요.

'정밀 수정'의 핵심 요구는 AI에게 내리는 모든 지시에 수정 범위와 경계를 명확히 지정하는 거예요. 'calculate\_discount 함수만 수정', '다른 부분은 건드리지 마', '기존 코드 스타일을 유지해'라고 말해 주세요. AI의 수정이 요청 범위를 넘어서면 즉시 남는 부분을 되돌리라고 요구하세요.

Karpathy는 diff(차이 비교) 예시로 이 점을 설명했어요. AI에게 장바구니의 할인 계산 함수에 '회원가' 로직을 추가하게 한다고 해 볼까요. 좋은 결과물은 이 함수 하나만 수정하고 몇 줄을 추가하는 것이고, 나쁜 결과물은 파일 전체를 다시 리팩터링해 버리는 거예요. 새 코드가 더 '우아해' 보일 수는 있지만, 새로운 버그를 끌어들이지 않았는지는 확신할 수 없죠. AI 프로그래밍에서 '더 좋아 보임'이 지시를 벗어나도 되는 이유가 되는 일은 결코 없어요.

나쁜 지시(범위가 모호함)

\begin{lstlisting}
이 프로젝트의 할인 계산 로직을 최적화해 줘
\end{lstlisting}

좋은 지시(범위를 정밀하게 지정)

\begin{lstlisting}
cart.py의 calculate\_discount 함수만 수정: 1. 기존 할인에 회원가 로직 추가 2. 회원은 추가로 10\% 할인 3. 함수 시그니처와 다른 함수는 수정 금지 4. 기존 변수명 스타일 유지
\end{lstlisting}

\booksection{원칙 4: Goal-Driven Execution(목표 기반 실행)}

이 원칙은 핵심 질문에 답해요. AI에게 '어떻게 할지'를 지시해야 할까요, '무엇을 할지'를 지시해야 할까요? Karpathy의 답은 후자예요. AI에게 '먼저 클래스를 만들고, 메서드를 추가하고, 테스트를 작성해'라고 말하는 대신, 'CSV 파일을 읽어 각 열의 고유값 개수를 세어 주는 도구가 필요해'라고 말하는 게 낫다는 거죠. 구현 경로는 AI 스스로 정하게 하고, 여러분은 결과 검수를 담당해요.

'목표 기반'의 장점은 두 배예요. AI 입장에서는 더 큰 자율 공간을 얻어 여러분의 중간 단계에 발목 잡히지 않고 최적의 구현 경로를 선택할 수 있어요. 여러분 입장에서는 관심을 '어떻게 하는가'에서 '결과가 맞는가'로 옮기게 되죠. 바로 지휘관이 해야 할 일이에요.

\begin{tcbwarn}{주의: 목표 기반 $\neq$ 방치하기}
AI가 실행한 뒤 엄격한 검수가 필요해요. 결과가 요구사항을 충족하나요? 경계 상황을 처리했나요? 눈에 띄는 과도한 설계는 없나요? 결과가 기대와 다르면 구체적인 피드백을 주세요. '다시 써'가 아니라 '여기서 문제는 빈 리스트를 입력하면 IndexError가 난다는 거예요. 빈 값 검사를 추가하도록 수정해 줘'처럼요.

\end{tcbwarn}

\begin{center}
\begin{tabularx}{\linewidth}{|X|X|X|}
\hline
\textbf{지시 스타일} & \textbf{예시} & \textbf{결과의 특징} \tabularnewline
\hline
명령형(어떻게 할지) & 'DataParser라는 Python 클래스를 먼저 만들고, parse 메서드를 추가하고, filename 매개변수를 받아 csv 모듈로 읽어…' & AI는 그저 타자수가 되고, 여러분이 정한 중간 단계가 최적 경로가 아닐 수도 있어요 \tabularnewline
목표형(무엇을 할지) & 'CSV 파일을 읽어 각 열의 고유값 개수를 세고 JSON 형식으로 출력하는 도구가 필요해요. 전체 구현을 보여 주세요.' & AI가 스스로 최적 경로를 선택하고, 여러분은 결과 검수를 담당해요 \tabularnewline
\hline
\end{tabularx}
\end{center}

\begin{center}\small\color{gray}명령형 지시와 목표형 지시 비교\end{center}

\booksection{원칙 5: Iterate and Verify(반복하고 검증하기)}

이 책이 Karpathy의 네 원칙 위에 더한 다섯 번째 원칙이에요. 핵심 믿음은 이것이에요. AI의 첫 결과물은 거의 언제나 최종본이 아니라는 거죠. 효과적인 AI 프로그래밍은 일회성 '지시-결과물'이 아니라 '생성-검증-수정'의 반복 순환이에요.

'반복 검증'은 수동적인 '안 되면 다시 해 보기'가 아니라 구조를 갖춘 검증 절차예요. AI가 결과물을 내놓을 때마다 세 층위로 검증해요.

\begin{itemize}
  \item 1층 · 기능 검증: 돌아가긴 하나요? 입력과 출력이 맞나요?
  \item 2층 · 경계 검증: 빈 입력, 지나치게 긴 입력, 비정상 입력을 넣으면 어떻게 되나요?
  \item 3층 · 아키텍처 검증: 코드 구조는 합리적인가요? 하드코딩된 부분은 없나요? 나중에 확장하기 쉬운가요?
\end{itemize}

목표 정의

AI가 코드 생성

기능 검증

통과

경계 검증

통과

아키텍처 검토

통과, 다음 단계로

미통과

구체적인 피드백 제공

다시 생성

\begin{center}\small\color{gray}그림 2: AI 프로그래밍의 반복 검증 순환\end{center}

이 순환의 관건은 '구체적인 피드백 주기' 단계예요. 검증을 통과하지 못했을 때 '틀렸어, 다시 해'라고 하지 말고 '이 함수는 빈 리스트를 입력하면 IndexError가 나요. 15번째 줄에 빈 값 검사를 추가해 줘'라고 말하세요. 구체적인 피드백은 AI가 정밀하게 고치게 하고, 모호한 피드백은 여러분의 의도를 다시 추측하게 만들 뿐이에요.

\booksection{Harness Engineering: Agent = Model + Harness}

2026년 중반, Harness Engineering이라는 개념이 AI 프로그래밍 커뮤니티에서 유행하기 시작했어요. 핵심 공식은 아주 간단해요.

에이전트의 능력 = 모델 능력 × 하네스 품질

모델 능력은 AI 두뇌의 지능 상한선이에요. Harness는 모델을 둘러싸고 구축한 '작업 환경'——컨텍스트 관리, 도구 통합, 흐름 오케스트레이션, 기억 시스템 등을 포함해요.

이 개념은 흔한 의문을 설명해 줘요. 똑같이 GPT-4나 Claude를 쓰면서 어떤 사람은 놀라운 프로젝트를 만들고 어떤 사람은 평범한 코드만 얻는 이유는 뭘까요? 차이는 모델이 아니라 Harness에 있어요. 좋은 Harness에는 명확한 프로젝트 규범 파일, 합리적인 도구 통합, 정돈된 대화 컨텍스트, 잘 갖춰진 검증 절차가 포함돼요. 모델이 지능 상한선을 정하고, Harness가 그 상한선을 얼마나 끌어낼 수 있는지를 정해요.

\begin{tcbinfo}{천재 인턴 비유}
모델을 천재 인턴이라고 상상해 보세요. IQ는 아주 높지만, 프로젝트 배경 문서를 주지 않고, 어떤 도구를 쓸지 알려 주지 않고, 기한과 검수 기준도 정해 주지 않으면 좋은 결과물을 만들지 못해요. Harness Engineering은 바로 '천재 인턴에게 작업 환경을 잘 마련해 주는 법'에 대한 엔지니어링이에요.

\end{tcbinfo}

이 책의 뒤 장들은 본질적으로 모두 좋은 Harness를 구축하는 법을 가르쳐 줘요. 제39장은 컨텍스트 관리(프로젝트 기억 시스템), 제40장은 요구사항 정의(PRD), 제43장은 대화 기술(Prompt 계층화), 제45장은 품질 보증(검증 절차)을 다루죠. 이것들을 합치면 바로 여러분의 Harness예요.

\begin{tcbpractice}{직접 해 보기: '생각의 삼각형'으로 다음 아이디어 뜯어보기}
늘 만들고 싶으면서 미뤄 둔 작은 것 하나를 골라 보세요(예: 디데이 웹페이지, 물 마시기 알림 스크립트). 3분만 들여 답해 보세요. (1) 무슨 문제를 해결하나요? (2) 만들면 어떤 모습인가요? (3) 어떤 제약(시간, 기술, 환경)이 있나요? 그리고 답을 한 문단으로 적어 보세요. 축하해요. 인생 첫 'Think Before Coding'을 완성했어요. 이 문단이 바로 합격점에 오른 고 품질 지시문이에요.

\end{tcbpractice}

\begin{tcbsummary}{이번 장 정리}
\begin{itemize}
  \item 5가지 원칙은 Karpathy의 70줄짜리 CLAUDE.md(Star 70K+)에서 왔고, 이 책은 다섯 번째 '반복 검증'을 보탰어요
  \item Think Before Coding: 말하기 전에 세 가지를 답하세요. 무엇을 하나? 결과는 어떤 모습? 어떤 제약이 있나?
  \item Simplicity First: AI의 가장 큰 함정은 '너무 많이 만들기'——할인 계산에 전자상거래 엔진은 필요 없어요
  \item Surgical Changes: 외과 수술처럼 정밀하게 수정하고, AI의 '지나가다 한번 손보기'는 거절해요
  \item Goal-Driven Execution: '어떻게'가 아니라 '무엇을'이라고 말하고, 경로 선택은 AI에게 맡겨요
  \item Iterate and Verify: 3층 검증(기능/경계/아키텍처) + 구체적 피드백으로 순환 반복해요
  \item Harness Engineering: 능력 = 모델 × Harness. AI를 '천재 인턴'처럼 대하고 환경을 관리해요
\end{itemize}

\end{tcbsummary}
